% syllabus/apstats-syllabus-sy2627.tex — AP Statistics course syllabus, Lynn English High School, SY 2026–27.
% Build: cd syllabus && pdflatex --miktex-enable-installer apstats-syllabus-sy2627.tex (twice)
% Facts come from the repo: roster-server/grade-config.js (weights, gates, quarters), SCHEDULE_HANDOFF.md
% (calendar, meeting days, exam date), APS_DOK_LADDER_SPEC.md (the daily flow), start-here.html (the promise).
\documentclass[10pt]{article}
\usepackage[letterpaper,margin=0.7in]{geometry}
\usepackage[T1]{fontenc}
\usepackage[utf8]{inputenc}
\usepackage{lmodern}
\usepackage{amsmath}
\usepackage{booktabs}
\usepackage{enumitem}
\usepackage{xcolor}
\usepackage[most]{tcolorbox}
\usepackage{hyperref}
\hypersetup{colorlinks=true, urlcolor=blue!60!black, linkcolor=black}
\setlength{\parindent}{0pt}
\setlength{\parskip}{5pt}
\definecolor{dayblue}{HTML}{2B5C8A}
\definecolor{callout}{HTML}{EAF1F8}
\definecolor{warm}{HTML}{FCEAE8}
\definecolor{warmred}{HTML}{C0392B}

\newcommand{\banner}[1]{%
  \vspace{4pt}\begin{tcolorbox}[enhanced, boxrule=0pt, arc=1pt, colback=dayblue, left=6pt, right=6pt, top=3pt, bottom=3pt]
  {\color{white}\bfseries\large #1}\end{tcolorbox}\vspace{-2pt}}

\pagestyle{empty}
\begin{document}

{\LARGE\bfseries AP Statistics}\hfill{\large Lynn English High School \textperiodcentered\ 2026--27}\\
{\large Mr.~Colson}\hfill{\small Period B (Mon/Tue/Thu/Fri) \textperiodcentered\ Period E (Mon/Wed/Fri)}\\
{\small Class site: \url{https://robjohncolson.github.io/apstats-live-worksheet/} \textperiodcentered\ AP Exam: Tuesday, May 11, 2027}
\vspace{2pt}\hrule

\banner{What this course is}
AP Statistics is a college-level course in collecting data, describing it, and using it to decide what is true.
It follows the College Board's Fall 2026 Course and Exam Description in five units, in this order:

\begin{enumerate}[leftmargin=1.6em, itemsep=1pt, topsep=2pt]
  \item Exploring One-Variable Data \& Collecting Data
  \item Probability, Random Variables \& Sampling Distributions
  \item Inference for Proportions \& Categorical Data
  \item Inference for Means
  \item Exploring Two-Variable Data (Regression)
\end{enumerate}
After the exam we keep going with ``Beyond the Exam'' topics that are worth knowing even though they are not tested.

\banner{My promise to you}
\begin{tcolorbox}[enhanced, colback=callout, colframe=dayblue, boxrule=0.6pt, arc=2pt, left=7pt, right=7pt, top=4pt, bottom=4pt]
\textbf{A grade should tell the truth about what you know.} One rough test day should not outweigh everything you can
actually do. This class is mastery-based and the door back is always open: fall behind, have a bad week, bomb a check,
and you can still go back, fix it, and show it again. Anyone here can earn full marks by finishing the loop: learn it, show
it, and where it is shaky, fix it and show it again. My job is to keep that path clear and keep it open.
\end{tcolorbox}

\banner{What a class period looks like}
\begin{enumerate}[leftmargin=1.6em, itemsep=1pt, topsep=2pt]
  \item \textbf{The problem goes on the board at the bell.} You get a paper sheet with one real statistics problem on it.
        Before anything is taught, you write a one-sentence \emph{first take}. No wrong answers yet.
  \item \textbf{The lesson video and follow-along.} On your Chromebook you open the day's lesson on the class site, watch the
        video, and fill in the follow-along worksheet as you go. The watching isn't the learning; the doing is.
  \item \textbf{Finish the problem.} With the video done and the rules printed on your sheet, you finish parts (a), (b), (c).
        Part (c) always asks you to \emph{decide something and defend it}: which claim the data support, which method is
        right, what a reader would get wrong. That is the thinking the AP exam rewards.
  \item \textbf{Turn the sheet in.} I read every one and hand it back. This sheet is the one piece of work each day that is
        graded by me, by hand, not by software.
\end{enumerate}
The lesson quiz and the flashcard round (Blooket) come after the video, in class when there is time, otherwise at home.

\banner{Tools you will use}
\begin{itemize}[leftmargin=1.4em, itemsep=1pt, topsep=2pt]
  \item \textbf{The Desk} on the class site: your home page. It shows today's lesson, what is due, your grade, and every resource.
  \item \textbf{Follow-along worksheets}: fill-in-the-blank while you watch, with reflection questions graded by an AI reader
        against a rubric. You can revise a worksheet any time, and a revision can only raise your score.
  \item \textbf{Lesson quizzes}: short, AP-style, with partial credit. A wrong answer is the most useful thing on the page.
  \item \textbf{Flashcards}: a quick Blooket on each lesson. Score 80\% or better to unlock the next lesson.
  \item \textbf{The TI-84 trainer}: calculator fluency is a skill you build. The trainer turns it into muscle memory.
  \item \textbf{Progress Checks}: in class, on paper, paced like the real exam, one per unit.
\end{itemize}

\banner{How your grade is computed}
Your quarter grade comes from two tracks, so that the grade rewards \emph{either} steady work \emph{or} demonstrated mastery,
whichever is stronger.

\begin{center}\small
\begin{tabular}{@{}llr@{}}
\toprule
\textbf{Track} & \textbf{What is in it} & \textbf{Weight inside the track} \\
\midrule
Work & Lessons (follow-along worksheets)   & 30\% \\
     & Lesson quizzes                      & 30\% \\
     & Posters (one per unit)              & 30\% \\
     & Flashcards (Blooket)                & 10\% \\
\midrule
Mastery & Progress Checks (in class, paper, exam-paced) & 100\% \\
\bottomrule
\end{tabular}
\end{center}

\begin{tcolorbox}[enhanced, colback=callout, colframe=dayblue, boxrule=0.6pt, arc=2pt, left=7pt, right=7pt, top=4pt, bottom=4pt]
\textbf{The rule.} If both tracks are at least 40\%, your quarter grade is the \textbf{higher} of the two.
If either track is below 40\%, the grade is the best of: 70\% of your Progress Check average, 70\% of your work average,
or the plain average of the two. Until the first Progress Check of a quarter exists, your grade is simply your work track.
\end{tcolorbox}

\begin{itemize}[leftmargin=1.4em, itemsep=1pt, topsep=2pt]
  \item \textbf{Due dates.} A lesson counts toward the work track once its class day has passed. Undone work that is past due
        counts as zero until you do it; the moment you do it, it counts. Nothing is ever locked.
  \item \textbf{Early finish bonus.} Finish a lesson (80\% or more of it) before its class day and you earn one bonus point on
        that lesson, up to five per quarter.
  \item \textbf{Progress Checks are required.} Skip one after it is due and your grade is capped until you take it. Walk in
        with the work done and it is upside, not a threat.
  \item \textbf{Revisions and appeals.} Worksheets can be revised at any time; AI grading only ever raises a score, and you can
        appeal a rubric score up to three times with your reasoning.
  \item \textbf{Schoology} mirrors these grades in five categories (Lesson, Quizzes, Posters, Blooket, Progress Check) and is
        updated from the class gradebook. The class gradebook on the Desk is the calculation of record; if the two ever
        disagree, ask me and I will fix it.
\end{itemize}

\banner{Calendar}
\begin{center}\small
\begin{tabular}{@{}llll@{}}
\toprule
\textbf{Quarter 1} & Sep 2 -- Nov 6 & \textbf{Quarter 3} & Jan 25 -- Apr 14 \\
\textbf{Quarter 2} & Nov 9 -- Jan 22 & \textbf{Quarter 4} & Apr 15 -- Jun 17 \\
\bottomrule
\end{tabular}
\end{center}
The AP Exam is \textbf{Tuesday, May 11, 2027}. Everything after it is Beyond-the-Exam material and review.

\banner{What I ask of you}
\begin{itemize}[leftmargin=1.4em, itemsep=1pt, topsep=2pt]
  \item Do the work honestly, and don't move on while something is broken.
  \item Bring a pencil every day; the daily problem is on paper. Bring your Chromebook charged; the lesson is on it.
  \item Use the AI tutor and the AI grader as coaches. The paper sheet you turn in is your own hand and your own thinking.
  \item Treat a weak quiz as information, not a verdict. Go back, shore it up, keep going.
  \item Show up for the Progress Check. It is the one thing on the calendar that cannot move.
\end{itemize}

\vspace{6pt}
\begin{tcolorbox}[enhanced, colback=warm, colframe=warmred, boxrule=0.6pt, arc=2pt, left=7pt, right=7pt, top=4pt, bottom=4pt]
\small\textbf{Questions?} Ask in class, message me through the Desk, or find me in my room before school.
\end{tcolorbox}

\end{document}
